\documentclass[aspectratio=169]{beamer}
% Paper is 16cm x 9cm. Compile with: latexmk -pdf -interaction=nonstopmode <file>.tex
% (two passes are required: the footline navigator and every `remember picture`
%  overlay read positions/total-frame counts from the .aux file).

\usetheme{default}
\usefonttheme{professionalfonts}

\usepackage[utf8]{inputenc}
\usepackage[T1]{fontenc}
\usepackage[<<BABEL_LANGUAGE>>]{babel}
\IfFileExists{montserrat.sty}{\usepackage{montserrat}}{\usepackage[scaled=0.95]{helvet}}
\usepackage{graphicx}
\usepackage{caption}
\usepackage{xcolor}
\usepackage{tikz}
\usetikzlibrary{shapes.geometric, arrows.meta}

% ── Palette ──────────────────────────────────────────────────────────────────
% Core (structural) colors
\definecolor{NoirBg}{HTML}{111317}      % canvas
\definecolor{NoirPanel}{HTML}{1A1F26}   % title band, zebra rows, panels
\definecolor{NoirText}{HTML}{F2EEE6}    % body text
\definecolor{NoirMuted}{HTML}{CFC8BB}   % secondary text, captions, borders
\definecolor{NoirAccent}{HTML}{D39A3C}  % amber: headings, bullets, progress
% Highlight colors (data emphasis, diagrams)
\definecolor{NoirCyan}{HTML}{00D4C8}    % "best"/positive values, code tokens
\definecolor{NoirPink}{HTML}{D81B9F}    % negative/alert values
\definecolor{NoirNavy}{HTML}{0A2540}    % fills for diagram shapes
\definecolor{NoirLink}{HTML}{00ACC1}    % rules, outlines, connectors
\colorlet{NoirFlow}{blue!55!white}     % diagram arrows / flow connectors

\setbeamercolor{background canvas}{bg=NoirBg}
\setbeamercolor{normal text}{fg=NoirText,bg=NoirBg}
\setbeamercolor{frametitle}{fg=NoirText,bg=NoirPanel}
\setbeamercolor{title}{fg=NoirAccent}
\setbeamercolor{subtitle}{fg=NoirText}
\setbeamercolor{structure}{fg=NoirAccent}
\setbeamercolor{itemize item}{fg=NoirAccent}
\setbeamercolor{itemize subitem}{fg=NoirAccent}
\setbeamercolor{caption name}{fg=NoirMuted}
\setbeamercolor{caption}{fg=NoirMuted}
\setbeamerfont{title}{series=\bfseries,size=\Large}
\setbeamerfont{frametitle}{series=\bfseries,size=\large}
\setbeamerfont{normal text}{size=\normalsize}
\renewcommand{\familydefault}{\sfdefault}

% ── Navigation: progress squares (bottom-right), no beamer nav symbols ──────
% One square per frame; amber = visited/current, dimmed = upcoming.
\setbeamertemplate{navigation symbols}{}
\setbeamertemplate{footline}{%
  \leavevmode%
  \hbox{%
    \begin{beamercolorbox}[wd=\paperwidth,ht=2.4ex,dp=1.1ex,leftskip=0pt,rightskip=0.8em]{author in head/foot}%
      \hfill%
      \begin{tikzpicture}[baseline=(current bounding box.center)]
        \pgfmathtruncatemacro{\CurrentFrame}{\insertframenumber}
        \pgfmathtruncatemacro{\TotalFrames}{max(1,\inserttotalframenumber)}
        \foreach \i in {1,...,\TotalFrames} {%
          \pgfmathsetmacro{\x}{(\i-1)*0.15}
          \ifnum\i>\CurrentFrame
            \fill[NoirMuted!35!NoirBg] (\x,0) rectangle ++(0.10,0.10);
          \else
            \fill[NoirAccent] (\x,0) rectangle ++(0.10,0.10);
          \fi
        }
      \end{tikzpicture}%
    \end{beamercolorbox}%
  }%
}
\setbeameroption{hide notes}

% ── Per-slide backgrounds ────────────────────────────────────────────────────
% Convention: every slide block is  \NoirBackground...  +  frame  +  \NoirBackgroundReset
% so a block can be inserted, moved or deleted without affecting its neighbours.
% A missing image renders a labelled placeholder instead of failing the build.
\newcommand{\NoirMissingImage}[3]{% #1 width, #2 height, #3 file name
  {\setlength{\fboxsep}{0pt}\colorbox{NoirPanel}{\parbox[c][#2][c]{#1}{\centering
    \color{NoirMuted}\ttfamily\scriptsize missing image: \detokenize{#3}}}}}
% Full-bleed image covering the whole slide (stretched to 16:9).
\newcommand{\NoirBackgroundImage}[1]{%
  \usebackgroundtemplate{%
    \IfFileExists{#1}%
      {\includegraphics[width=\paperwidth,height=\paperheight]{#1}}%
      {\NoirMissingImage{\paperwidth}{\paperheight}{#1}}}}
% Pure black slide with an image banner hung just below the title band.
\newcommand{\NoirBackgroundBanner}[1]{%
  \usebackgroundtemplate{%
    \begin{tikzpicture}[remember picture, overlay]
      \fill[black] (current page.south west) rectangle (current page.north east);
      \node[anchor=north, inner sep=0pt] at ([yshift=-0.13\paperheight] current page.north) {%
        \IfFileExists{#1}%
          {\includegraphics[width=\paperwidth]{#1}}%
          {\NoirMissingImage{\paperwidth}{0.55\paperheight}{#1}}};
    \end{tikzpicture}}}
% Pure black slide (used for diagrams, tables and text-only cards).
\newcommand{\NoirBackgroundVoid}{%
  \usebackgroundtemplate{%
    \begin{tikzpicture}[remember picture, overlay]
      \fill[black] (current page.south west) rectangle (current page.north east);
    \end{tikzpicture}}}
% Back to the default NoirBg canvas. NOTE: \usebackgroundtemplate{} is NOT a reset:
% it installs an empty template and the slide renders WHITE. Restore the default.
\newcommand{\NoirBackgroundReset}{\setbeamertemplate{background canvas}[default]}

% ── Metadata ─────────────────────────────────────────────────────────────────
\title[<<SHORT_TITLE>>]{<<TITLE>>}
\subtitle{<<SUBTITLE>>}
\author{<<AUTHOR>>}
\institute{<<INSTITUTE_OR_COURSE>>}
\date{\today}

\begin{document}

% ── Slides go here ────────────────────────────────────────────────────────────
% Each slide is delimited by a banner comment:  % === Slide: <name> ===

% === Slide: Title ===
\begin{frame}
\centering
\begin{tikzpicture}[x=0.95cm, y=0.95cm,
    every node/.style={inner sep=0pt, outer sep=0pt},
    bighex/.style={regular polygon, regular polygon sides=6,
        draw=NoirCyan, line width=1.2pt, fill=NoirNavy, minimum size=2.6cm,
        text=NoirCyan, font=\small\bfseries, align=center, inner sep=3pt,
        text width=1.5cm},  % keep the core keyword to 1–2 short words
    sathex/.style={regular polygon, regular polygon sides=6,
        fill=NoirCyan, draw=none, minimum size=1.4cm,
        text=NoirBg, font=\tiny\bfseries, align=center, inner sep=2pt},
    dimhex/.style={regular polygon, regular polygon sides=6,
        draw=NoirCyan!50, line width=0.7pt, fill=NoirNavy!80, minimum size=1.1cm,
        text=NoirText, font=\tiny, align=center, inner sep=2pt},
    dot/.style={circle, fill=NoirPink, minimum size=0.34cm},
    netline/.style={draw=NoirCyan!45, line width=0.7pt},
    dotline/.style={draw=NoirPink!65, line width=0.6pt}
]
% Invisible canvas frame (keeps the bounding box stable): 14.7 x 8.6 units = 13.97 x 8.17cm,
% i.e. exactly the 16:9 text area. Keep every element inside it.
\path (0,0) rectangle (14.7,8.6);

% ── Left: kicker + title + rule ──────────────────────────────────────────────
\node[anchor=north west, font=\LARGE\bfseries, text=NoirAccent]
    at (0.20, 8.50) {\insertshorttitle};
\node[anchor=north west, font=\large\bfseries, text=NoirText,
      text width=6.5cm, align=left]
    at (0.20, 7.75) {\inserttitle};
\node[anchor=north west, font=\small, text=NoirMuted,
      text width=6.5cm, align=left]
    at (0.20, 5.05) {\insertsubtitle};
\draw[NoirLink, line width=0.7pt] (0.20, 5.30) -- (6.80, 5.30);

% ── Right: decorative hexagon network (relabel or replace with \includegraphics)
\node[bighex] (core) at (11.60, 4.50) {<<CORE_KEYWORD>>};
\node[sathex] (n1) at (11.60, 6.80) {};
\node[dimhex] (n2) at (14.00, 5.40) {};
\node[dimhex] (n3) at (9.20, 5.60) {};
\node[dimhex] (n4) at (9.10, 3.20) {};
\node[sathex, minimum size=1.1cm] (n5) at (11.60, 2.30) {};
\node[dot] (d1) at (12.40, 5.90) {};
\node[dot] (d2) at (12.90, 4.10) {};
\node[dot] (d3) at (10.40, 3.40) {};
\foreach \a/\b in {core/n1, core/n2, core/n3, core/n4, core/n5, n1/n2, n3/n1, n4/n3, n4/n5}
    \draw[netline] (\a) -- (\b);
\foreach \a/\b in {d1/n1, d1/n2, d2/n2, d2/n5, d3/n4, d3/n5}
    \draw[dotline] (\a) -- (\b);

% ── Bottom: byline box ───────────────────────────────────────────────────────
\node[draw=NoirLink, line width=0.9pt, rounded corners=4pt,
      font=\small, text=NoirText, text width=12.5cm, align=center, inner sep=7pt]
    at (7.50, 0.55) {\insertauthor, \insertinstitute};
\end{tikzpicture}
\end{frame}

% ── Insert new slides above this line ────────────────────────────────────────
\end{document}
